%% SCRAP: archive/operations/jenkins-pipeline-guide
%% SOURCE: docs/working/archive/operations/jenkins-pipeline-guide.adoc
%% STATUS: OBSOLETE
%% FITS: none
%% EDITORIAL: lifted — prose rewritten to press voice

\section{Jenkins Developer Pipeline Guide (2025-10-30)}

This document described the StarForth Jenkins developer pipeline as it existed at
version 1.0 (2025-10-30). The pipeline is no longer current; it is preserved as a
historical operations record.

\subsection{Pipeline Overview}

The developer baseline pipeline (\texttt{Jenkinsfile}) ran in approximately two hours
and covered 17~stages:

\begin{enumerate}
  \item Workspace cleanup and preparation.
  \item Build gauntlet: DEBUG, STANDARD, FASTEST, and FAST configurations in parallel.
  \item Smoke test.
  \item Comprehensive test suite (936 tests).
  \item Benchmark gauntlet: quick (1\,M iterations), full suite, and stack, arithmetic,
        and logic stress tests (10\,M iterations each).
  \item Extreme stress tests: deep recursion, maximum stack depth, memory allocation
        patterns, long-running stability, and nested loops.
  \item Thermal and performance monitoring under sustained load.
  \item Memory leak detection (Valgrind).
  \item Profile-guided optimisation build.
  \item PGO performance comparison.
  \item Edge-case testing: division by zero, stack underflow, stack overflow, invalid
        memory access.
  \item Performance regression check against baseline CSV.
  \item Code quality: compiler warnings and lines-of-code analysis.
  \item Formal verification: Isabelle theory build (audit mode), refinement status
        check, code annotation validation, refinement report generation.
  \item Documentation build: Doxygen API reference, \LaTeX{} conversion.
  \item Package build: Debian (\texttt{.deb}) and Red Hat (\texttt{.rpm}) via
        \texttt{fpm}.
  \item Test report generation.
\end{enumerate}

\subsection{Planned Variants}

Three derived pipelines were documented but not implemented at the time of archival:
a Test pipeline (stages 1--8, 13--17, approximately one hour), a QA pipeline (stages
1--7, 13--17 plus artifact signing, approximately 45 minutes), and a Production pipeline
(stages 1--4, 13--17 plus security scanning and repository upload, approximately
15 minutes).

%% TODO(bob): confirm current CI/CD pipeline status and whether Jenkinsfiles were
%% superseded by GitHub Actions or another system.
